\documentclass{article}
\usepackage[T1]{fontenc}
\usepackage{lmodern}
\usepackage{graphicx}
\usepackage{amsmath,amssymb}
\usepackage[a4paper,margin=2cm]{geometry}
\usepackage[absolute,overlay]{textpos}
\setlength{\TPHorizModule}{1mm}
\setlength{\TPVertModule}{1mm}
\usepackage[indonesian]{babel}
\usepackage{hyperref}
\setlength{\emergencystretch}{2em}
% unit: Halaman 82

\begin{document}

% === Halaman 82 (llm) ===

\textbf{Step 3:} Simpan salinan state awal ke dalam larik x

\begin{verbatim}
for (i = 0; i < 16; ++i)
    x[i] = in[i];
\end{verbatim}

\textbf{Step 4:} Lakukan 20 putaran (= 10 putaran ganda), satu putaran ganda terdiri dari 10 \textit{column round} dan 10 \textit{row round}.

Double-round 1
$\rightarrow$ Column Round
$\rightarrow$ Row Round

Double-round 2
$\rightarrow$ Column Round
$\rightarrow$ Row Round

$\dots$

Double-round 10
$\rightarrow$ Column Round
$\rightarrow$ Row Round

\end{document}
